\chapter{Sécurité, autorisation et traçabilité}
\label{ch:securite}

\section{Authentification JWT}

L'authentification repose sur des jetons JSON Web Token signés en HMAC-SHA-512. À la
connexion, le serveur émet un jeton d'accès valable une heure et un jeton de
rafraîchissement valable sept jours. Le jeton d'accès porte l'identifiant du compte, son
adresse électronique, son type et son rôle.

Les mots de passe ne sont jamais stockés en clair : ils sont hachés par BCrypt avec un
facteur de coût de~12, soit $2^{12}$ itérations. Ce paramètre rend une attaque par
dictionnaire coûteuse tout en gardant une vérification de l'ordre de quelques centaines de
millisecondes, acceptable pour une opération de connexion.

Le mécanisme comporte un choix de conception qui mérite d'être explicité, car il va à
l'encontre de l'usage le plus répandu.

\paragraph{Le rôle inscrit dans le jeton n'est pas utilisé pour autoriser.} À chaque
requête, \code{JwtAuthenticationFilter} vérifie la signature, puis recharge l'utilisateur
depuis la base de données et reconstruit ses habilitations à partir de l'état courant. Le
rôle présent dans la charge utile n'est conservé qu'à titre informatif pour le client.

Ce choix a un coût --- une lecture supplémentaire par requête --- et un bénéfice
déterminant : un jeton ne peut pas survivre à la révocation des droits qu'il transporte. Un
membre mis à la retraite pendant la durée de vie de son jeton voit son accès refusé à la
requête suivante, sans qu'il soit nécessaire de tenir une liste de jetons révoqués.

\section{Autorisation RBAC}

Le contrôle d'accès est fondé sur les rôles et s'applique à deux niveaux complémentaires.

Le premier niveau, déclaré dans \code{SecurityConfig}, distingue les ressources publiques
des ressources protégées : les endpoints d'authentification, le préfixe \code{/public/**},
la documentation de l'API et les lectures des publications, événements, équipes et
partenaires sont ouverts ; tout le reste exige une authentification.

Le second niveau, activé par \code{@EnableMethodSecurity}, place une annotation
\code{@PreAuthorize} sur chaque méthode ou classe de contrôleur. La règle est ainsi lisible
à côté du code qu'elle protège. Ce niveau produit la matrice détaillée à la section
suivante.

Un troisième cas, plus rare, relève du service et non de l'annotation : lorsque
l'autorisation dépend de l'identité de la cible et pas seulement du rôle de l'appelant. La
modification d'un profil en est l'exemple --- un membre peut modifier le sien, la direction
peut modifier celui d'un autre --- et la vérification est portée par
\code{MemberService}.

\section{Matrice des permissions}

Le tableau~\ref{tab:matrice} présente la matrice complète, relevée des annotations du code
et vérifiée par appels réels sur l'application en fonctionnement.

\begin{landscape}
% En mode paysage la hauteur de composition retombe à la largeur du bloc de texte :
% l'interligne des 36 lignes est resserré pour que la matrice tienne sur une seule
% page, sans page blanche intercalaire ni débordement dans la marge basse.
\begin{table}[H]
\centering
\captionsetup{skip=3pt}
\caption{Matrice des permissions par rôle applicatif ($\bullet$ = autorisé)}
\label{tab:matrice}
\footnotesize
\renewcommand{\arraystretch}{0.90}
\begin{tabular}{@{}p{7.6cm}ccccccc@{}}
\toprule
\textbf{Opération} & \rotatebox{0}{\textbf{ADMIN}} & \textbf{DIR} & \textbf{V-DIR} & \textbf{CHEF} & \textbf{PERM} & \textbf{ASSO} & \textbf{DOCT} \\
\midrule
\multicolumn{8}{@{}l}{\itshape Consultation générale} \\
Portail public, publications, événements, équipes & $\bullet$ & $\bullet$ & $\bullet$ & $\bullet$ & $\bullet$ & $\bullet$ & $\bullet$ \\
Liste et fiches des membres & $\bullet$ & $\bullet$ & $\bullet$ & $\bullet$ & $\bullet$ & $\bullet$ & $\bullet$ \\
Calendrier, notifications, discussions, recherche & $\bullet$ & $\bullet$ & $\bullet$ & $\bullet$ & $\bullet$ & $\bullet$ & $\bullet$ \\
Modifier son propre profil & $\bullet$ & $\bullet$ & $\bullet$ & $\bullet$ & $\bullet$ & $\bullet$ & $\bullet$ \\
\midrule
\multicolumn{8}{@{}l}{\itshape Activité scientifique} \\
Ajouter / modifier une publication & $\bullet$ & $\bullet$ & $\bullet$ & $\bullet$ & $\bullet$ & $\bullet$ & $\bullet$ \\
Supprimer une publication & $\bullet$ & $\bullet$ & $\bullet$ & & & & \\
Créer / modifier un événement ou une réunion & $\bullet$ & $\bullet$ & $\bullet$ & $\bullet$ & & & \\
Déposer un procès-verbal & $\bullet$ & $\bullet$ & $\bullet$ & $\bullet$ & & & \\
Supprimer un événement, une réunion & $\bullet$ & $\bullet$ & $\bullet$ & & & & \\
S'inscrire à un événement & $\bullet$ & $\bullet$ & $\bullet$ & $\bullet$ & $\bullet$ & $\bullet$ & $\bullet$ \\
\midrule
\multicolumn{8}{@{}l}{\itshape Modules internes (fermés aux doctorants --- CdC §2.2.C)} \\
Consulter documents, matériel, procès-verbaux & $\bullet$ & $\bullet$ & $\bullet$ & $\bullet$ & $\bullet$ & $\bullet$ & \\
Déposer un document, créer une version & $\bullet$ & $\bullet$ & $\bullet$ & $\bullet$ & $\bullet$ & $\bullet$ & \\
Supprimer un document & $\bullet$ & $\bullet$ & $\bullet$ & & & & \\
Soumettre une demande de matériel & $\bullet$ & $\bullet$ & $\bullet$ & $\bullet$ & $\bullet$ & $\bullet$ & \\
\midrule
\multicolumn{8}{@{}l}{\itshape Gestion des ressources} \\
Gérer l'inventaire, distribuer, enregistrer un retour & $\bullet$ & $\bullet$ & $\bullet$ & & & & \\
Valider ou rejeter une demande de matériel & & $\bullet$ & $\bullet$ & & & & \\
Gérer partenaires et conventions & $\bullet$ & $\bullet$ & $\bullet$ & & & & \\
\midrule
\multicolumn{8}{@{}l}{\itshape Gouvernance et membres} \\
Modifier le profil d'un autre membre & $\bullet$ & $\bullet$ & $\bullet$ & & & & \\
Changer le statut d'un membre & $\bullet$ & $\bullet$ & $\bullet$ & & & & \\
Transférer un membre entre équipes & $\bullet$ & $\bullet$ & $\bullet$ & & & & \\
Archiver / restaurer un membre, départ, retour & $\bullet$ & $\bullet$ & & & & & \\
Créer, modifier, archiver une équipe & $\bullet$ & $\bullet$ & & & & & \\
Attribuer ou clôturer un mandat & $\bullet$ & $\bullet$ & & & & & \\
Consulter les demandes d'adhésion & $\bullet$ & $\bullet$ & $\bullet$ & & & & \\
\textbf{Statuer sur une demande d'adhésion} & & $\bullet$ & & & & & \\
\midrule
\multicolumn{8}{@{}l}{\itshape Pilotage et administration} \\
Générer et télécharger le rapport annuel & $\bullet$ & $\bullet$ & $\bullet$ & & & & \\
Consulter le journal d'audit & $\bullet$ & $\bullet$ & $\bullet$ & & & & \\
Gérer les comptes et les rôles techniques & $\bullet$ & & & & & & \\
Modifier l'identité du laboratoire & $\bullet$ & & & & & & \\
\bottomrule
\end{tabular}
\end{table}
\end{landscape}

Trois asymétries de cette matrice ne sont pas des erreurs mais des traductions littérales
du cahier des charges.

La décision sur une demande d'adhésion appartient au seul directeur (§6 : « uniquement par
le directeur en mandat »). L'administrateur système, pourtant le plus privilégié
techniquement, n'y a pas accès : c'est une décision institutionnelle, non technique.

Symétriquement, la validation d'une demande de matériel appartient au directeur et au
vice-directeur, mais pas à l'administrateur système, pour la même raison.

Enfin, l'administrateur système ne gère pas les équipes ni les mandats au même titre que le
directeur --- il les partage avec lui --- mais il est seul à pouvoir créer un compte ou
modifier l'identité du laboratoire.

La figure~\ref{fig:sidebarPhd} illustre l'effet de cette matrice sur l'interface : la
navigation d'un doctorant ne comporte que les six modules qui lui sont ouverts. Les
modules internes ne sont pas grisés, ils sont absents.

\begin{figure}[H]
\centering
\includegraphics[width=\textwidth]{\figpath/screens/24-sidebar-doctorant.png}
\caption{Navigation réduite d'un compte doctorant --- les modules internes sont absents}
\label{fig:sidebarPhd}
\end{figure}

\section{Gestion des statuts}

Les cinq statuts d'appartenance sont portés par l'énumération \code{MemberStatus}. Un
changement de statut est toujours historisé dans \code{MemberStatusHistory} avec sa date
d'effet et son motif, et déclenche systématiquement la resynchronisation du rôle.

Deux statuts produisent un effet de bord sur la sécurité : le passage à
\code{RETIRED} ou \code{FORMER} clôt les mandats en cours en plus de geler le compte.
L'archivage d'un membre applique la même règle, pour la même raison : quelle que soit la
manière dont un membre cesse d'exercer, il ne conserve aucune responsabilité de gouvernance
ouverte. Cette clôture explicite garantit que l'historique de gouvernance reste cohérent
--- on ne trouve pas de mandat ouvert détenu par un membre parti ou archivé.

\section{Synchronisation rôle / mandat}

Le mécanisme a été décrit en §\ref{sec:srm} et illustré par la figure~\ref{fig:seqMandat}.
Du point de vue de la sécurité, sa propriété essentielle est la suivante : il n'existe
aucun chemin par lequel un rôle de gouvernance puisse être obtenu sans mandat
correspondant, à l'exception de la dérogation administrative explicite
(\code{ADMIN\_OVERRIDE}), elle-même historisée et signalée dans l'interface.

La réconciliation exécutée au démarrage constitue un filet de sécurité supplémentaire.
Elle garantit qu'un compte dont le rôle aurait dérivé --- par import, par intervention
directe en base ou par antériorité --- est ramené à la valeur que son statut et ses mandats
impliquent, avec une trace portant la cause \code{RECONCILIATION}. Sur le système en
exploitation, cette procédure ne corrige plus rien et signale simplement que tous les
comptes sont conformes.

\section{Protection des comptes retraités et anciens}

Le gel repose sur une seule expression, évaluée à la construction des détails
d'authentification :

\begin{lstlisting}[language=Java,caption={Règle de gel des comptes (\code{CustomUserDetails})},label={lst:frozen}]
// Un membre RETIRED ou FORMER est gele, quelle que soit la valeur brute
// de users.enabled (CdC section 8 : "compte gele" / "compte desactive").
this.enabled = user.isEnabled() && !isFrozenStatus(memberStatus);
\end{lstlisting}

Cette expression est évaluée à chaque requête, puisque les détails sont reconstruits à
chaque passage du filtre. Un compte gelé ne peut donc ni se connecter ni continuer à
utiliser un jeton déjà émis. Son historique demeure en revanche intégralement consultable
par la direction, conformément à la section~8 du cahier des charges.

Un compte technique dépourvu de profil de membre --- l'administrateur système --- n'est pas
concerné par ce gel : l'absence de statut n'est pas assimilée à un statut gelé.

\section{Journal d'audit}

Toute action sensible est consignée : création, modification, suppression, restauration,
approbation, rejet, connexion, échec de connexion, déconnexion, réinitialisation de mot de
passe, changement de rôle, changement de statut, dépôt et téléchargement de fichier.

Chaque entrée conserve l'identité de l'acteur, son adresse électronique, l'action, le type
et l'identifiant de l'entité concernée, un libellé, l'adresse IP et l'horodatage. L'usage
de la propagation \code{REQUIRES\_NEW} garantit que la trace subsiste même si l'opération
auditée est annulée.

Le journal est consultable par l'administrateur système et la direction, avec filtrage par
entité et par acteur.

\section{Gestion des secrets}

Aucun identifiant réel ne figure dans le code source suivi en gestion de version. Les
paramètres de connexion à la base, les identifiants de messagerie, la clé de signature
des jetons et le répertoire de stockage sont lus depuis des variables d'environnement. Un
fichier local, exclu du dépôt, les fournit en développement ; un fichier d'exemple ne
contenant que des valeurs fictives documente les clés attendues.

\paragraph{La clé de signature est obligatoire.} Le fichier de configuration de base
s'applique à \emph{tous} les profils : une valeur par défaut qui y figurerait serait
active en production. La clé y est donc déclarée sans valeur de repli, et son absence
arrête le démarrage. Trois contrôles sont appliqués avant que le premier jeton ne soit
émis : la clé doit être présente, être un base64 décodable, et peser au moins 64~octets
une fois décodée. Ce dernier seuil tient à la façon dont l'algorithme est choisi : il
n'est jamais nommé dans le code, il est déduit de la clé. Avec jjwt~0.12.6,
\code{Keys.hmacShaKeyFor} refuse une clé de moins de 256~bits, retient HS256 de 256 à
383~bits, HS384 de 384 à 511~bits, et HS512 à partir de 512~bits --- soit exactement
64~octets. Une clé de 32 à 63~octets ne provoquerait donc aucune erreur : elle ferait
signer les jetons en HS256 ou HS384 au lieu du HS512 annoncé plus haut. Le profil de test fournit
sa propre clé, sans valeur de secret et identifiée comme telle dans son contenu même.

Ce point a été corrigé tardivement : il est décrit parmi les défauts relevés
en~§\ref{sec:defauts}.

Cette externalisation a été mise en place en fin de projet, après le constat que des
identifiants de base de données figuraient dans un fichier de configuration versionné. La
remédiation a comporté deux volets : la suppression des valeurs du code suivi, et la
rotation effective du mot de passe concerné --- vérifiée par test de connexion : le nouvel
identifiant est accepté, l'ancien est rejeté. L'historique de gestion de version conserve
encore l'ancienne valeur, désormais sans effet ; la procédure de purge est documentée dans
le dépôt et reste à la discrétion du propriétaire.

Cet épisode constitue un enseignement méthodologique du projet : une configuration de
développement est un fichier de code comme un autre, et doit être traitée comme tel dès le
premier commit.

\paragraph{Le cas particulier du compte d'amorçage.} Une base de données neuve ne contient
aucun compte, et toutes les routes d'administration sont protégées : le tout premier
administrateur système doit donc être créé par le programme lui-même. Ce privilège est le
plus étendu du système et il est encadré par quatre garde-fous. Le composant d'amorçage
est restreint au profil \code{dev}, de sorte qu'un déploiement de production n'obtient
jamais un compte qu'il n'a pas explicitement créé. Le mot de passe est lu dans une variable
d'environnement dépourvue de valeur de repli : en son absence, rien n'est créé plutôt
qu'un compte à secret par défaut. Un secret de moins de douze caractères est refusé.
Enfin le compte ainsi créé porte l'indicateur de changement de mot de passe obligatoire,
comme tout compte provisionné par l'administration ou par l'acceptation d'une demande
d'adhésion. Un compte déjà présent n'est jamais réécrit : la rotation du secret d'une
instance en service relève de l'exploitation et s'effectue depuis l'écran de profil.
Ces quatre propriétés sont verrouillées par la suite \code{AdminAccountInitializerTest}
(§\ref{sec:defauts}).
